\section{The corrected best response}
\begin{proposition}[Derivative, concavity, and reduced best response]\label{prop:br}
Holding $\theta_j$ and $\theta_k$ fixed, direct differentiation of Eq.~\eqref{eq:profit} yields
\begin{equation}\label{eq:derivative}
\frac{\partial\pi_j}{\partial i_j}
=-\frac K2+
\frac{(1-\rho)r I_1-v(\theta_j-\theta_k)I_2
 +(1-\delta)^2\delta v(1-2\theta_j+\theta_k)
 -(2i_j-i_k)L}{2\tau}.
\end{equation}
Moreover,
\begin{equation}\label{eq:hessian}
\frac{\partial^2\pi_j}{\partial i_j^2}=-\frac{L}{\tau}
=-\frac{1+2(1-\delta)^3}{3\tau}<0.
\end{equation}
Hence the unique maximizer of the reduced one-dimensional objective is
\begin{equation}\label{eq:BR}
\boxed{
 i_j(i_k)=\frac{i_k}{2}+
\frac{(1-\rho)r I_1-v(\theta_j-\theta_k)I_2
 +(1-\delta)^2\delta v(1-2\theta_j+\theta_k)-\tau K}{2L}.}
\end{equation}
\end{proposition}

\begin{proof}
Differentiate the risky-account block in Eq.~\eqref{eq:profit}. The integrals multiplying the own rate are
\[
\int_0^\delta(1-\lambda)\dd\lambda=I_1,\qquad
\int_0^\delta\lambda(1-\lambda)\dd\lambda=I_2,
\]
and
\[
\int_0^\delta(1-\lambda)^2\dd\lambda=I_3.
\]
The risky block therefore contributes
\[
-\frac{I_1}{2}+
\frac{(1-\rho)rI_1-v(\theta_j-\theta_k)I_2-(2i_j-i_k)I_3}{2\tau}.
\]
Differentiating the liquid block contributes
\[
-\frac{(1-\delta)^2}{2}+
\frac{(1-\delta)^2\delta v(1-2\theta_j+\theta_k)
 -(2i_j-i_k)(1-\delta)^3}{2\tau}.
\]
Adding the two expressions gives Eq.~\eqref{eq:derivative}. Differentiating once more gives Eq.~\eqref{eq:hessian}. Strict concavity then makes the first-order condition sufficient, and solving it yields Eq.~\eqref{eq:BR}.
\end{proof}

The difference from the source's displayed Eq.~(8) is not a normalization. In particular, the term $(1-\rho)r I_1$ is generated by the interaction between the risky-deposit investment margin and the rate-induced change in market share. The source's displayed Hessian also omits part of the curvature. The corrected Hessian is negative throughout the reduced domain, which strengthens rather than weakens the concavity property of that one-dimensional problem.

\subsection{An exact regression point}
The discrepancy can be exhibited without relying on numerical optimization. Set
\begin{equation}\label{eq:counterpars}
\delta=\frac12,\quad \rho=0,\quad v=1,\quad r=\frac12,\quad
\tau=1,\quad \theta_A=\theta_B=0.
\end{equation}
The source formulas imply $i_A=i_B=-1$. At Eq.~\eqref{eq:counterpars}, $I_1=3/8$, $K=5/8$, and $L=5/12$. Equation~\eqref{eq:derivative} gives
\[
\left.\frac{\partial\pi_A}{\partial i_A}\right|_{i_A=i_B=-1}=\frac{5}{96}>0,
\]
so the candidate is not stationary. The corrected symmetric rate is $-3/4$.

This is also a valid same-regime deviation. If bank $A$ moves from $-1$ to $-99/100$ while bank $B$ remains at $-1$, direct evaluation of Eq.~\eqref{eq:profit} gives
\[
\pi_A(-99/100,-1)-\pi_A(-1,-1)=\frac{1}{2000}>0.
\]
Throughout $\lambda\in[0,1/2]$, the risky share is
\[
 x_R(\lambda)=\frac{101}{200}-\frac{\lambda}{200}\in\left[\frac{201}{400},\frac{101}{200}\right],
\]
and $x_L=201/400$. Thus the profitable deviation remains strictly inside the same affine market-share regime.
